% fig_gallery: one small schematic per model family (Sec. IV). Few nodes, few words.
% Only the useful or partly useful models get a schematic (Vivian, 2026-10-07): 02, 11, 16, 17, 03. No Potts, no 01 on its own.
% Ising: vermillion/blue nodes (spin states). Vector spins: arrows colored by angle (cyclic 'twilight' map, sampled). Field: orange.
\definecolor{tw25}{HTML}{647DBC}\definecolor{tw40}{HTML}{6066B6}\definecolor{tw50}{HTML}{5F55AF}
\definecolor{tw60}{HTML}{5E43A5}\definecolor{tw70}{HTML}{5B3298}\definecolor{tw85}{HTML}{531D7A}\definecolor{tw100}{HTML}{421257}
\definecolor{pottsA}{HTML}{009E73}
\begin{figure*}[t]
\centering
\resizebox{0.86\textwidth}{!}{%
\begin{tikzpicture}[font=\sffamily\small,
  ag/.style={circle,draw=inkc,line width=0.5pt,minimum size=5mm,inner sep=0pt},
  rl/.style={-{Stealth[length=1.5mm]},draw=inkc!70,line width=0.5pt},
  ttl/.style={anchor=west,font=\sffamily\small\bfseries},
  lab/.style={font=\sffamily\footnotesize,text=inkc}]
% (a) Ising, M02/01
\node[ttl] at (-0.3,3.3) {(a) M02 kinetic Ising};
\node[ag,fill=spinup!85] (i1) at (0.3,2.3) {}; \node[ag,fill=spinup!85] (i2) at (1.4,2.6) {};
\node[ag,fill=spindn!85] (i3) at (2.4,1.9) {}; \node[ag,fill=spinup!85] (i4) at (0.6,1.0) {};
\node[ag,fill=spindn!85] (i5) at (1.8,0.7) {};
\draw[rl] (i1)--(i2); \draw[rl] (i2)--(i3); \draw[rl] (i4)--(i5); \draw[rl] (i5)--(i3);
\node[lab,align=center] at (1.3,-0.45) {flip at own call;\\links act at read-out};
% (b) vector spins, M11
\begin{scope}[xshift=3.6cm]
\node[ttl] at (-0.3,3.3) {(b) M11 vector spins};
\draw[-{Stealth[length=2.6mm]},draw=fieldc,line width=2.4pt] (2.65,0.5) -- (3.05,2.75);
\node[lab,text=fieldc,anchor=west] at (2.95,2.9) {goal};
\foreach \x/\y/\a/\c in {0.3/2.1/75/tw40, 1.3/2.25/105/tw85, 0.5/1.0/85/tw60, 1.6/1.3/65/tw25, 1.1/0.3/95/tw70, 2.1/2.0/80/tw50} {
  \fill[inkc!25] (\x,\y) circle (0.08);
  \draw[-{Stealth[length=2mm]},draw=\c,line width=1.6pt] (\x,\y) -- ++(\a:0.7);}
\node[lab,align=center] at (1.4,-0.45) {aligned to the field,\\weakly to each other};
\end{scope}
% (c) Langevin, M16
\begin{scope}[xshift=7.6cm]
\node[ttl] at (-0.3,3.3) {(c) M16 Langevin};
\draw[->,draw=nullc] (0,0.5) -- (3.0,0.5) node[lab,anchor=north east] at (3.0,0.45) {calls};
\draw[->,draw=nullc] (0,0.5) -- (0,2.9);
\draw[draw=inkc,line width=0.9pt,domain=0:2.9,smooth,samples=40,variable=\t] plot (\t,{0.5+1.9*exp(-\t/1.6)});
\draw[draw=spindn,line width=1.1pt,domain=0:2.9,smooth,samples=40,variable=\t] plot (\t,{0.5+1.9*exp(-\t/0.18)});
\node[lab,anchor=west] at (0.9,1.85) {well};
\node[lab,anchor=west,text=spindn] at (0.25,0.85) {kick};
\node[lab,align=center] at (1.5,-0.45) {two rates,\\ratio $\approx$15};
\end{scope}
% (d) collective modes, M17
\begin{scope}[xshift=11.3cm]
\node[ttl] at (-0.3,3.3) {(d) M17 modes};
\draw[->,draw=nullc] (0,0.5) -- (3.1,0.5) node[lab,anchor=north east] at (3.1,0.45) {eigenvalue};
\fill[nullc!45,domain=0.15:1.45,smooth,variable=\x] (0.15,0.5) -- plot (\x,{0.5+2.2*sqrt(max(0,(\x-0.15)*(1.45-\x)))}) -- (1.45,0.5) -- cycle;
\draw[dashed,draw=inkc] (1.55,0.5) -- (1.55,2.6) node[lab,anchor=south] {edge};
\fill[spindn] (2.6,0.5) rectangle (2.72,1.6);
\node[lab,anchor=south,text=spindn] at (2.66,1.6) {mode};
\node[lab,align=center] at (1.5,-0.45) {noise bulk;\\one mode above};
\end{scope}
% (e) branching, M03
\begin{scope}[xshift=15.0cm]
\node[ttl] at (-0.3,3.3) {(e) M03 branching};
\node[ag,fill=pottsA!80,minimum size=4mm] (r) at (1.4,2.5) {};
\node[ag,fill=pottsA!50,minimum size=4mm] (c1) at (0.8,1.5) {};
\node[ag,fill=pottsA!50,minimum size=4mm] (c2) at (2.0,1.5) {};
\node[ag,fill=pottsA!30,minimum size=4mm] (g1) at (2.0,0.6) {};
\draw[rl] (r)--(c1); \draw[rl] (r)--(c2); \draw[rl] (c2)--(g1);
\node[lab,align=center] at (1.4,-0.45) {$\hat R\approx0.2$:\\trees die out};
\end{scope}
\end{tikzpicture}}
\caption{The useful model families, from the most to the least useful. (a)~Kinetic Ising with read-out coupling (M02): each agent is a two-state spin (two colors) that updates only at its own calls; links act at the read-out call. (b)~Vector spins (M11): each agent's content is an arrow (color = angle); the goal is a strong field (orange), so the arrows align with it and only weakly with each other. (c)~Langevin relaxation (M16): a read kick decays about 15 times faster than the agent's own well. (d)~Collective modes (M17): a mode counts only above the calibrated noise edge. (e)~Branching (M03): idea trees die out, $\hat R\approx0.2$.}
\label{fig:gallery}
\end{figure*}
